\documentclass[addpoints,12pt]{exam}

\usepackage{geometry}
 \geometry{
 a4paper,
 total={170mm,260mm},
 left=20mm,
 top=15mm,
 }

\usepackage{amsmath,amssymb,amsthm}
\DeclareFontEncoding{T1}
\fontencoding{T1}
\usepackage[T2A]{fontenc} 
\usepackage[utf8]{inputenc}
\usepackage[russian]{babel}
\usepackage{graphics,graphicx}
\usepackage{indentfirst}

\usepackage{verbatim}
\usepackage{fancyvrb}  
\usepackage{booktabs}
\usepackage{multirow}

\usepackage{listings}
\usepackage{color}

\renewcommand{\solutiontitle}{\noindent\textbf{Решение:}\par\noindent}

\definecolor{gray}{gray}{0.5}

\lstset{basicstyle=\ttfamily,
    language=Matlab,    
    keepspaces=true,
    extendedchars=\true,
    identifierstyle={ \bf \color{black} },
    showstringspaces=false,
    numbers=left,%
    numbersep=7pt,  
    emph=[1]{case,switch,otherwise,nonlocal,as,yield,zip,print},    
    frame=single,    
    xleftmargin=0.3cm,
    frame=l,framesep=4pt,framerule=0.5pt,
    numberstyle={\scriptsize \color{gray}}
}


\pagestyle{headandfoot}
\runningheadrule
\firstpageheader{ИнтМатПак}{MATLAB 1}{10 февраля 2026}
\runningheader{ИнтМатПак}
{MATLAB 1, Стр. \thepage\ из \numpages}
{10 февраля 2026}
\firstpagefooter{}{}{}
\runningfooter{}{}{}
\pointname{} 
\boxedpoints
%\pointsinmargin
\pointsinrightmargin


%\header{MATLAB 1}
%}
%{
%	\hfill 29.05.2023
%}

% \footer{Заведующий кафедрой математических методов в экономике}
% {}
% {
%     \hfill проф. Гераськин М. И.
% }


\begin{document}
%\begin{center}
%\fbox{\fbox{\parbox{5.5in}{\centering
%Правильный.}}}
%\end{center}
\vspace{0.1in}
\makebox[\textwidth]{Ф. И. О., группа: \enspace\hrulefill}
\vspace{0.1in}

\begin{questions}

\question[1]
Какой последний элемент будет содержаться в матрице B в результате выполнения этого кода?
\begin{lstlisting}
A = reshape(1:64,8,8);
B = A(3:6, 3:6); 
\end{lstlisting}
\begin{checkboxes}
\CorrectChoice 46;
\choice 45; 
\choice 38;
\choice 37; 
\choice все ответы неверны.
\end{checkboxes}

\question[1] 
Напишите код, который извлекает в переменную C каждый второй столбец матрицы A из предыдущего задания:
\begin{solutionorlines}[1cm]
\end{solutionorlines}

\question[1] Кинетическая энергия углового движения твердого тела определяется матричным выражением $ T = \boldsymbol{\omega}^T \boldsymbol{J} \boldsymbol{\omega} / 2$, где $\boldsymbol{\omega}$ - координатный столбец угловой скорости тела, $\boldsymbol{J}$ - матрица (3x3) тензора инерции тела. В сессии MATLAB определены переменные 
\begin{lstlisting}
J = [1,2,3]; % диагональные элементы диагональной матрицы тензора инерции 
w = [0.5, 0.2, 0.3]; % компоненты угловой скорости
\end{lstlisting}
Выберите верный код для вычисления кинетической энергии:
\begin{checkboxes}
\choice  \ttfamily{T = 0.5*w*J'*w'};
\choice  \ttfamily{T = 0.5*w'*J*w};
\CorrectChoice  \ttfamily{T = 0.5*w*diag(J)*w'};
\choice  \ttfamily{T = 0.5*w*diag(J)*w'};
\end{checkboxes}

\question[1] Чему будет равна переменная b?
\begin{lstlisting}
a = [7,3,6,8,2,4];
b = reshape(a,2,3);
\end{lstlisting}
\begin{checkboxes}
\choice  [7, 3, 8; 6, 2, 4]
\CorrectChoice [7, 6, 2; 3, 8, 4]
\choice  [7, 3; 8, 6; 2, 4]
\choice  [7, 6; 2, 3; 8, 4]
\end{checkboxes}

\question[1] Есть матрица temperatures размером 30×24, представляющая температуру (в градусах Цельсия) каждый час в течение 30 дней (строка — день, столбец — час). Напишите код, который определяет среднюю температуру для каждого дня (результат — вектор-столбец 30×1).
\begin{solutionorlines}[2cm]
\end{solutionorlines}
 
\question[1] Чему будет равна переменная i в результате выполнения кода?
\begin{lstlisting}
a = [7,3,6,4];
[~,i] = sort(a);
\end{lstlisting}
\begin{checkboxes}
\CorrectChoice \texttt{[2,4,3,1]}
\choice \texttt{[1,3,4,2]}
\choice \texttt{[3,4,6,7]}
\choice \texttt{[7,6,4,3]}
\end{checkboxes}

\question[1] Чему будет равна переменная a?
\begin{lstlisting}
a = sum([7,6,3,9].*[0,1,0,1])
\end{lstlisting}    
\begin{checkboxes}
\choice 7;
\choice 9;
\CorrectChoice 15;
\choice выведется сообщение об ошибке.
\end{checkboxes}

\question[1] Координаты двух точек в пространстве заданы матрицами-строками \texttt{A = [1,2,0]; B = [3,1,1]}. Напишите код, который определяет расстояние между этими двумя точками.
\begin{solutionorlines}[1cm]
\end{solutionorlines}

\question[1] 
Дана матрица X размером 50x50, заполненная нормально-распределенными случайными числами (мат. ожидание равно 20, стандартное отклонение равно 10): 
\begin{lstlisting}
X = randn(50, 50) * 10 + 20; 
\end{lstlisting}
Напишите код, который извлекает <<периметр>> матрицы X (первая и последняя строки, первый и последний столбцы) в линейный вектор \texttt{perimeter\_vec}, идущий по часовой стрелке, начиная с элемента (1,1).
\begin{solutionorlines}[2cm]
    \texttt{perimeter\_vec} = 
\end{solutionorlines}

\question[1] 
Напишите код, который в последнем столбце матрицы из предыдущего задания (\texttt{X(:, end)}) все значения, превышающие 30, увеличивает в два раза (должна изменится матрица X).
\begin{solutionorlines}[2cm]
\end{solutionorlines}

\end{questions}

\hqword{Вопрос:}
\hpword{Баллы:}
\hsword{Ответ:}
\htword{Сумма}

\vfill

\begin{center}
\gradetable[h][questions]
\end{center}

% Отлично: 12,13; Хорошо: 8,9,10; Удовлетворительно 5,6,7

\end{document}
